% ====================================================================
%  徐森林、薛春华《数学分析》（第一册）
%  \S3.4  L'Hospital 法则
%  源：原书印刷页 198-206（PDF 物理页 204-212）
%  底本：PaddleOCR-VL-1.6 + 高清影印比对校订
%  规范：全西文半角标点；单一落盘出口 _texout.py
% ====================================================================
\section{L'Hospital 法则}\label{sec:3.4}

本节介绍的 L'Hospital(洛必塔)法则是求 $\frac{0}{0}$, $\frac{\infty}{\infty}$, $0\cdot\infty$, $0^{0}$, $1^{\infty}$, $\infty^{0}$, $\infty-\infty$ 各种不定型极限的非常有效的方法.

\begin{theorem}[$\frac{0}{0}$ 不定型极限, L'Hospital 法则]\label{thm:3.4.1}
设 $f$, $g$ 在区间 $(x_0,b)$ 上可导, 且
\[
\lim_{x\to x_0^{+}}f(x)=0=\lim_{x\to x_0^{+}}g(x),\quad g'(x)\neq0,\quad \forall x\in(x_0,b).
\]
如果 $\lim\limits_{x\to x_0^{+}}\frac{f'(x)}{g'(x)}=A$(实数, $\pm\infty$, $\infty$), 那么
\[
\lim_{x\to x_0^{+}}\frac{f(x)}{g(x)}=\lim_{x\to x_0^{+}}\frac{f'(x)}{g'(x)}=A.
\]
\end{theorem}

\begin{proof}
延拓 $f$, $g$ 到 $x_0$, s.t. $f(x_0)=g(x_0)=0$, 则 $f$, $g$ 在 $[x_0,b]$ 上连续, 且 $\forall x\in(x_0,b)$, $g'(x)\neq0$. 根据 Cauchy 中值定理知, $\exists\xi\in(x_0,x)$, s.t.
\[
\frac{f(x)}{g(x)}=\frac{f(x)-f(x_0)}{g(x)-g(x_0)}=\frac{f'(\xi)}{g'(\xi)}.
\]
令 $x\to x_0^{+}$, 此时, $\xi\to x_0^{+}$, 所以,
\[
\lim_{x\to x_0^{+}}\frac{f(x)}{g(x)}=\lim_{x\to x_0^{+}}\frac{f'(\xi)}{g'(\xi)}=\lim_{x\to x_0^{+}}\frac{f'(x)}{g'(x)}.
\]
\end{proof}

\begin{remark}\label{rem:3.4.1}
当 $x\to x_0^{-}$, $x\to x_0$ 时, 类似定理~\ref{thm:3.4.1}, 有相同的结论.
\end{remark}

\begin{theoremp}[$\frac{0}{0}$ 不定型极限, L'Hospital 法则]\label{thm:3.4.1p}
设 $f$, $g$ 在 $(a,+\infty)$ 上可导, 且
\[
\lim_{x\to+\infty}f(x)=0=\lim_{x\to+\infty}g(x),\quad g'(x)\neq0,\quad \forall x\in(a,+\infty).
\]
如果 $\lim\limits_{x\to+\infty}\frac{f'(x)}{g'(x)}=A$(实数, $\pm\infty$, $\infty$), 那么
\[
\lim_{x\to+\infty}\frac{f(x)}{g(x)}=\lim_{x\to+\infty}\frac{f'(x)}{g'(x)}=A.
\]
\end{theoremp}

\begin{proof}
令 $x=\frac{1}{t}$, 当 $x\to+\infty$ 时, $t\to0^{+}$, 则
\[
\lim_{x\to+\infty}f(x)=0\Longleftrightarrow\lim_{t\to0^{+}}f\left(\frac{1}{t}\right)=0,
\]
\[
\lim_{x\to+\infty}g(x)=0\Longleftrightarrow\lim_{t\to0^{+}}f\left(\frac{1}{t}\right)=0.
\]
于是,
\[
\begin{aligned}
\lim_{x\to+\infty}\frac{f(x)}{g(x)}&=\lim_{t\to0^{+}}\frac{f\left(\frac{1}{t}\right)}{g\left(\frac{1}{t}\right)}=\lim_{t\to0^{+}}\frac{f'\left(\frac{1}{t}\right)\cdot\frac{-1}{t^{2}}}{g'\left(\frac{1}{t}\right)\cdot\frac{-1}{t^{2}}}=\lim_{t\to0^{+}}\frac{f'\left(\frac{1}{t}\right)}{g'\left(\frac{1}{t}\right)}\\
&=\lim_{x\to+\infty}\frac{f'(x)}{g'(x)}.
\end{aligned}
\]
\end{proof}

\begin{remarkp}\label{rem:3.4.1p}
当 $x\to-\infty$, $x\to\infty$ 时, 类似定理~\ref{thm:3.4.1p} 有相同的结论.
\end{remarkp}

\begin{theorempp}[$\frac{0}{0}$ 不定型极限, 推广 L'Hospital 法则]\label{thm:3.4.1pp}
设当 $x\to x_0^{+}$ 时, $\frac{f}{g}$, $\frac{f'}{g'}$, $\cdots$, $\frac{f^{(n-1)}}{g^{(n-1)}}$ 均为 $\frac{0}{0}$ 型, 且 $g^{(i)}(x)\neq0$, $i=1,\cdots,n$. 如果 $\lim\limits_{x\to x_0^{+}}\frac{f^{(n)}(x)}{g^{(n)}(x)}=A$(实数, $\pm\infty$, $\infty$), 则
\[
\lim_{x\to x_0^{+}}\frac{f(x)}{g(x)}=\lim_{x\to x_0^{+}}\frac{f'(x)}{g'(x)}=\cdots=\lim_{x\to x_0^{+}}\frac{f^{(n)}(x)}{g^{(n)}(x)}=A.
\]
\end{theorempp}

\begin{proof}
由题设将定理~\ref{thm:3.4.1} 应用于 $f^{(n-1)}(x)$, $g^{(n-1)}(x)$ 推出
\[
\lim_{x\to x_0^{+}}\frac{f^{(n-1)}(x)}{g^{(n-1)}(x)}=\lim_{x\to x_0^{+}}\frac{f^{(n)}(x)}{g^{(n)}(x)}.
\]
依此类推得到
\[
\lim_{x\to x_0^{+}}\frac{f(x)}{g(x)}=\lim_{x\to x_0^{+}}\frac{f'(x)}{g'(x)}=\cdots=\lim_{x\to x_0^{+}}\frac{f^{(n-1)}(x)}{g^{(n-1)}(x)}=\lim_{x\to x_0^{+}}\frac{f^{(n)}(x)}{g^{(n)}(x)}.
\]
\end{proof}

\begin{remarkpp}\label{rem:3.4.1pp}
当 $x\to x_0^{-}$, $x\to x_0$, $x\to\pm\infty$, $x\to\infty$ 时, 类似定理~\ref{thm:3.4.1pp} 有相同的结论.
\end{remarkpp}

\begin{theorem}[$\frac{\cdot}{\infty}$ 不定型极限, L'Hospital 法则]\label{thm:3.4.2}
设 $f$, $g$ 都在 $(x_0,b)$ 上可导, 且
\[
\lim_{x\to x_0^{+}}g(x)=\infty,\quad g'(x)\neq0,\quad \forall x\in(x_0,b).
\]
如果 $\lim\limits_{x\to x_0^{+}}\frac{f'(x)}{g'(x)}=A$(实数, $\pm\infty$, $\infty$), 那么
\[
\lim_{x\to x_0^{+}}\frac{f(x)}{g(x)}=\lim_{x\to x_0^{+}}\frac{f'(x)}{g'(x)}=A.
\]
\end{theorem}

\begin{proof}[证法 1]
只证 $\lim\limits_{x\to x_0^{+}}\frac{f'(x)}{g'(x)}=A\in\mathbb{R}$ 的情形, 其他情形 $(A=\pm\infty,\infty)$ 类似证明.

$\forall\varepsilon>0$, $\exists b_1\in(x_0,b)$, s.t. 当 $x\in(x_0,b_1)$ 时,
\[
A-\frac{\varepsilon}{2}<\frac{f'(x)}{g'(x)}<A+\frac{\varepsilon}{2}.
\]
因此, 根据 Cauchy 中值定理, $\exists\xi\in(x,b_1)$, s.t.
\[
\begin{aligned}
A-\frac{\varepsilon}{2}<\left(\frac{f(x)}{g(x)}-\frac{f(b_1)}{g(x)}\right)\left(1-\frac{g(b_1)}{g(x)}\right)^{-1}&=\frac{f(x)-f(b_1)}{g(x)-g(b_1)}\\
&=\frac{f'(\xi)}{g'(\xi)}<A+\frac{\varepsilon}{2},\\
\left(1-\frac{g(b_1)}{g(x)}\right)\left(A-\frac{\varepsilon}{2}\right)+\frac{f(b_1)}{g(x)}&<\frac{f(x)}{g(x)}\\
&<\left(1-\frac{g(b_1)}{g(x)}\right)\left(A+\frac{\varepsilon}{2}\right)+\frac{f(b_1)}{g(x)}.
\end{aligned}
\]
又因为 $\lim\limits_{x\to x_0}g(x)=\infty$, 故
\[
\lim_{x\to x_0^{+}}\left[\left(1-\frac{g(b_1)}{g(x)}\right)\left(A-\frac{\varepsilon}{2}\right)+\frac{f(b_1)}{g(x)}\right]=A-\frac{\varepsilon}{2},
\]
\[
\lim_{x\to x_0^{+}}\left[\left(1-\frac{g(b_1)}{g(x)}\right)\left(A+\frac{\varepsilon}{2}\right)+\frac{f(b_1)}{g(x)}\right]=A+\frac{\varepsilon}{2}.
\]
由此知, $\exists b_2\in(x_0,b_1)$, s.t. 当 $x\in(x_0,b_2)$ 时,
\[
A-\varepsilon<\frac{f(x)}{g(x)}<A+\varepsilon.
\]
这就证明了
\[
\lim_{x\to x_0^{+}}\frac{f(x)}{g(x)}=A=\lim_{x\to x_0^{+}}\frac{f'(x)}{g'(x)}.
\]
\end{proof}

\begin{proof}[证法 2]
在不等式
\[
\begin{aligned}
&\left(1-\frac{g(b_1)}{g(x)}\right)\left(A-\frac{\varepsilon}{2}\right)+\frac{f(b_1)}{g(x)}<\frac{f(x)}{g(x)}\\
&<\left(1-\frac{g(b_1)}{g(x)}\right)\left(A+\frac{\varepsilon}{2}\right)+\frac{f(b_1)}{g(x)}
\end{aligned}
\]
两边取上、下极限得到
\[
A-\frac{\varepsilon}{2}\leqslant\varliminf_{x\to x_0^{+}}\frac{f(x)}{g(x)}\leqslant\varlimsup_{x\to x_0^{+}}\frac{f(x)}{g(x)}\leqslant A+\frac{\varepsilon}{2},
\]
再令 $\varepsilon\to0^{+}$, 有
\[
A\leqslant\varliminf_{x\to x_0^{+}}\frac{f(x)}{g(x)}\leqslant\varlimsup_{x\to x_0^{+}}\frac{f(x)}{g(x)}\leqslant A,
\]
\[
\lim_{x\to x_0}\frac{f(x)}{g(x)}=\lim_{x\to x_0^{+}}\frac{f(x)}{g(x)}=\varlimsup_{x\to x_0^{+}}\frac{f(x)}{g(x)}=A=\lim_{x\to x_0}\frac{f'(x)}{g'(x)}.
\]
\end{proof}

注意: 定理~\ref{thm:3.4.2} 中, 当 $x\to x_0^{+}$ 时, $\frac{\cdot}{\infty}$ 的分子可以趋于 $\infty$, 也可以不趋于 $\infty$.

关于 $\frac{\cdot}{\infty}$ 不定型极限, 类似定理~\ref{thm:3.4.1p}、定理~\ref{thm:3.4.1pp}、注~\ref{rem:3.4.1}、注~\ref{rem:3.4.1p}、注~\ref{rem:3.4.1pp} 可得到定理 3.4.2$'$、定理 3.4.2$''$、注 3.4.2$'$、注 3.4.2$''$, 不再一一赘述.

\begin{example}\label{ex:3.4.1}
在例~\ref{ex:2.2.13} 中, 应用 L'Hospital 法则得到
\[
\lim_{x\to+\infty}\frac{\ln x}{x^{k}}=\lim_{x\to+\infty}\frac{1}{x}\cdot\frac{1}{kx^{k-1}}=\frac{1}{k}\lim_{x\to+\infty}\frac{1}{x^{k}}=0,\quad k>0.
\]
\[
\begin{aligned}
\lim_{x\to+\infty}\frac{x^{k}}{a^{x}}&=\lim_{x\to+\infty}\frac{kx^{k-1}}{a^{x}\ln a}=\lim_{x\to+\infty}\frac{k(k-1)x^{k-2}}{a^{x}(\ln a)^{2}}=\cdots\\
&=k(k-1)\cdots(k-n+1)\lim_{x\to+\infty}\frac{x^{k-n}}{a^{x}(\ln a)^{n}}=0,
\end{aligned}
\]
其中, $n-1<k\leqslant n$, $a>1$.

进而, 当 $p(x)=a_0x^{n}+\cdots+a_{n-1}x+a_{n}$ 为多项式时,
\[
\begin{aligned}
\lim_{x\to+\infty}\frac{p(x)}{a^{x}}&=\lim_{x\to+\infty}\frac{p'(x)}{a^{x}\ln a}=\lim_{x\to+\infty}\frac{p''(x)}{a^{x}(\ln a)^{2}}=\cdots\\
&=\lim_{x\to+\infty}\frac{a_0n!}{a^{x}(\ln a)^{n}}=0,\quad a>0.
\end{aligned}
\]
此外, 还有
\[
\lim_{x\to+\infty}\frac{a^{x}}{x^{x}}=\lim_{x\to+\infty}\mathrm{e}^{x(\ln a-\ln x)}=0,\quad a>0.
\]
得到这一等式时, 根本不需要用 L'Hospital 法则. 实际上, 用 L'Hospital 法则不能求得极限.
\end{example}

\begin{example}\label{ex:3.4.2}
\[
\lim_{x\to0}\frac{x-\sin x}{\frac{x^{3}}{3!}}=\lim_{x\to0}\frac{1-\cos x}{\frac{x^{2}}{2!}}=\lim_{x\to0}\frac{\sin x}{x}=1.
\]
必须提醒读者注意: 加减不能用等价代换, 请检查下述计算哪一步是错误的.
\[
\lim_{x\to0}\frac{x-\sin x}{\frac{x^{3}}{3!}}=\lim_{x\to0}\frac{x-x}{\frac{x^{3}}{3!}}=\lim_{x\to0}0=0.
\]
\end{example}

在计算极限时, 会遇到“$0\cdot\infty$”、“$\infty-\infty$”、“$0^{0}$”、“$\infty^{0}$”、“$1^{\infty}$”等不定型. 只要经过适当的变换, 最终都可以化为“$\frac{0}{0}$”或“$\frac{\infty}{\infty}$”型的情形, 再应用 L'Hospital 法则去解决.

\begin{example}\label{ex:3.4.3}
求极限 $\lim\limits_{x\to0^{+}}x^{\alpha}\ln x$, 其中 $\alpha>0$.
\end{example}

\begin{solution}
\[
\begin{aligned}
\lim_{x\to0^{+}}x^{\alpha}\ln x&\xlongequal{0\cdot\infty}\lim_{x\to0^{+}}\frac{\ln x}{\left(\frac{1}{x}\right)^{\alpha}}\xlongequal{\frac{\infty}{\infty}}\lim_{x\to0^{+}}\frac{\frac{1}{x}}{\alpha x^{-\alpha-1}}\\
&=-\frac{1}{\alpha}\lim_{x\to0^{+}}x^{\alpha}=0.
\end{aligned}
\]
\end{solution}

\begin{example}\label{ex:3.4.4}
求极限:
(1) $\lim\limits_{x\to1}\left(\frac{1}{x-1}-\frac{1}{\ln x}\right)$; \quad (2) $\lim\limits_{x\to\pi/2}(\sec x-\tan x)$.
\end{example}

\begin{solution}
这都是“$\infty-\infty$”型, 但可化为 $\frac{0}{0}$ 型.

(1)
\[
\begin{aligned}
\lim_{x\to1}\left(\frac{1}{x-1}-\frac{1}{\ln x}\right)&=\lim_{x\to1}\frac{\ln x-(x-1)}{(x-1)\ln x}=\lim_{x\to1}\frac{\frac{1}{x}-1}{\frac{x-1}{x}+\ln x}\\
&=\lim_{x\to1}\frac{1-x}{x-1+x\ln x}=\lim_{x\to1}\frac{-1}{1+1+\ln x}=-\frac{1}{2}.
\end{aligned}
\]

(2)
\[
\lim_{x\to\pi/2}(\sec x-\tan x)=\lim_{x\to\pi/2}\frac{1-\sin x}{\cos x}=\lim_{x\to\pi/2}\frac{-\cos x}{-\sin x}=0.
\]
\end{solution}

\begin{example}\label{ex:3.4.5}
求极限:
(1) $\lim\limits_{x\to0^{+}}(\sin x)^{x}$; \quad (2) $\lim\limits_{x\to0^{+}}(\sin x)^{\frac{k}{1+\ln x}}$ ($k$ 为常数);
(3) $\lim\limits_{x\to0}(\cos x)^{\frac{1}{x^{2}}}$; \quad (4) $\lim\limits_{x\to\frac{\pi}{2}^{-}}(\tan x)^{\frac{\pi}{2}-x}$.
\end{example}

\begin{solution}
(1)
\[
\begin{aligned}
\lim_{x\to0^{+}}(\sin x)^{x}&\xlongequal{0^{0}}\exp\left(\lim_{x\to0^{+}}x\ln\sin x\right)\xlongequal{0\cdot\infty}\exp\left(\lim_{x\to0^{+}}\frac{\ln\sin x}{x^{-1}}\right)\\
&=\exp\left(\lim_{x\to0^{+}}\frac{1}{-x^{-2}}\cdot\frac{\cos x}{\sin x}\right)=\exp\left(\lim_{x\to0^{+}}\frac{-x}{\sin x}\cdot x\cos x\right)=\mathrm{e}^{0}=1.
\end{aligned}
\]

(2)
\[
\begin{aligned}
\lim_{x\to0^{+}}(\sin x)^{\frac{k}{1+\ln x}}&\xlongequal{0^{0}}\exp\left(\lim_{x\to0^{+}}\frac{k\ln\sin x}{1+\ln x}\right)=\exp^{k}\left(\lim_{x\to0^{+}}\frac{\frac{\cos x}{\sin x}}{x^{-1}}\right)\\
&=\exp^{k}\left(\lim_{x\to0^{+}}\cos x\cdot\frac{x}{\sin x}\right)=\mathrm{e}^{k}.
\end{aligned}
\]

(3)
\[
\lim_{x\to0}(\cos x)^{\frac{1}{x^{2}}}\xlongequal{1^{\infty}}\exp\left(\lim_{x\to0}\frac{\ln\cos x}{x^{2}}\right)=\exp\left(\lim_{x\to0}\frac{-\tan x}{2x}\right)=\mathrm{e}^{-\frac{1}{2}}.
\]

(4)
\[
\begin{aligned}
\lim_{x\to\frac{\pi}{2}^{-}}(\tan x)^{\frac{\pi}{2}-x}&\xlongequal{t=\frac{\pi}{2}-x}\lim_{t\to0^{+}}\left(\frac{\cos t}{\sin t}\right)^{t}=\frac{\lim\limits_{t\to0^{+}}(\cos t)^{t}}{\lim\limits_{t\to0^{+}}(\sin t)^{t}}\\
&=\frac{1}{1}=1.
\end{aligned}
\]
\end{solution}

应用 L'Hospital 法则, 可重新计算例~\ref{ex:2.2.12}(3) 及例~\ref{ex:2.2.13}(3) 中的极限.

\begin{example}\label{ex:3.4.6}
(1) 求 $\lim\limits_{x\to+\infty}\left(1+\frac{1}{n}-\frac{1}{n^{2}}\right)^{n}$;
(2) 求 $\lim\limits_{x\to0^{+}}\left(\frac{a_1^{x}+\cdots+a_m^{x}}{m}\right)^{\frac{1}{x}}$.
\end{example}

\begin{solution}
(1)
\[
\begin{aligned}
\lim_{x\to+\infty}\left(1+\frac{1}{n}-\frac{1}{n^{2}}\right)^{n}&\xlongequal{1^{\infty}}\lim_{x\to0}(1+x-x^{2})^{\frac{1}{x}}\\
&=\exp\left(\lim_{x\to0}\frac{\ln(1+x-x^{2})}{x}\right)=\exp\left(\lim_{x\to0}\frac{1-2x}{1+x-x^{2}}\right)=\mathrm{e}^{1}=\mathrm{e}.
\end{aligned}
\]

(2)
\[
\begin{aligned}
\lim_{x\to0^{+}}\left(\frac{a_1^{x}+\cdots+a_m^{x}}{m}\right)^{\frac{1}{x}}&\xlongequal{1^{\infty}}\exp\left(\lim_{x\to0^{+}}\frac{\ln(a_1^{x}+\cdots+a_m^{x})-\ln m}{x}\right)\\
&=\exp\left(\lim_{x\to0^{+}}\frac{a_1^{x}\ln a_1+\cdots+a_m^{x}\ln a_m}{a_1^{x}+\cdots+a_m^{x}}\right)\\
&=\exp\left(\frac{\ln a_1+\cdots+\ln a_m}{m}\right)=\mathrm{e}^{\ln\sqrt[m]{a_1\cdots a_m}}=\sqrt[m]{a_1\cdots a_m}.
\end{aligned}
\]
\end{solution}

注意: 不能在数列形式下直接用 L'Hospital 法则, 因为对离散变量 $n\in\mathbb{N}$ 是无法求导数的.

\begin{example}\label{ex:3.4.7}
设
\[
f(x)=\begin{cases}\frac{g(x)}{x}, & x\neq0,\\ 0, & x=0,\end{cases}
\]
且 $g(0)=g'(0)=0$, $g''(0)=3$. 试求 $f'(0)$.
\end{example}

\begin{solution}
\[
\begin{aligned}
f'(0)&=\lim_{x\to0}\frac{f(x)-f(0)}{x-0}=\lim_{x\to0}\frac{\frac{g(x)}{x}-0}{x-0}=\lim_{x\to0}\frac{g(x)}{x^{2}}\\
&\xlongequal{\frac{0}{0}}\lim_{x\to0}\frac{g'(x)}{2x}=\frac{1}{2}\lim_{x\to0}\frac{g'(x)-g'(0)}{x-0}=\frac{1}{2}g''(0)=\frac{3}{2}.
\end{aligned}
\]
如果用两次 L'Hospital 法则, 得到
\[
f'(0)=\lim_{x\to0}\frac{g(x)}{x^{2}}=\lim_{x\to0}\frac{g'(x)}{2x}=\lim_{x\to0}\frac{g''(x)}{2}=\frac{1}{2}g''(0)=\frac{3}{2}.
\]
错在何处?
\end{solution}

\begin{example}\label{ex:3.4.8}
设 $f$ 在点 $x$ 处 2 阶可导. 证明:
\[
f''(x)=\lim_{h\to0}\frac{f(x+h)-f(x-h)-2f(x)}{h^{2}}.
\]
\end{example}

\begin{proof}
因为 $f$ 在 $x$ 点 2 阶可导, 故 $\exists\delta>0$, s.t. $f$ 在 $(x-\delta,x+\delta)$ 中 1 阶可导. 于是,
\[
\begin{aligned}
&\lim_{h\to0}\frac{f(x+h)+f(x-h)-2f(x)}{h^{2}}\xlongequal[\text{L'Hospital 法则}]{\frac{0}{0}}\lim_{h\to0}\frac{f'(x+h)-f'(x-h)}{2h}\\
&=\frac{1}{2}\lim_{h\to0}\left[\frac{f'(x+h)-f'(x)}{h}+\frac{f'(x-h)-f'(x)}{-h}\right]=\frac{1}{2}(f''(x)+f''(x))=f''(x).
\end{aligned}
\]
\end{proof}

\begin{example}\label{ex:3.4.9}
下面两种计算哪个是正确的? 如果有错误, 请指出错在何处?
\[
\lim_{x\to+\infty}\frac{2x-\cos x}{2x+\cos x}=\lim_{x\to+\infty}\frac{2-\frac{\cos x}{x}}{2+\frac{\cos x}{x}}=\frac{2-0}{2+0}=1.
\]
\[
\begin{aligned}
\lim_{x\to+\infty}\frac{2x-\cos x}{2x+\cos x}&\xlongequal{\frac{\infty}{\infty}\text{ 型}}\lim_{x\to+\infty}\frac{2+\sin x}{2-\sin x}=\lim_{x\to+\infty}\frac{\cos x}{-\cos x}\\
&=\lim_{x\to+\infty}(-1)=-1.
\end{aligned}
\]
\end{example}

\begin{example}\label{ex:3.4.10}
设 $f$ 在 $(a,+\infty)$ 上可导, 并且 $\lim\limits_{x\to+\infty}[f(x)+f'(x)]=0$. 证明: $\lim\limits_{x\to+\infty}f(x)=0$.
\end{example}

\begin{proof}[证法 1]
由 $\frac{\cdot}{\infty}$ 型 L'Hospital 法则得
\[
\begin{aligned}
\lim_{x\to+\infty}f(x)&=\lim_{x\to+\infty}\frac{\mathrm{e}^{x}f(x)}{\mathrm{e}^{x}}\xlongequal{\frac{\cdot}{\infty}}\lim_{x\to+\infty}\frac{\mathrm{e}^{x}[f(x)+f'(x)]}{\mathrm{e}^{x}}\\
&=\lim_{x\to+\infty}[f(x)+f'(x)]=0.
\end{aligned}
\]
\end{proof}

\begin{proof}[证法 2]
因为 $\lim\limits_{x\to+\infty}[f(x)+f'(x)]=0$, 所以, $\exists\Delta_1>a$, s.t. 当 $x>\Delta_1$ 时, 有
\[
|f(x)+f'(x)|<\frac{\varepsilon}{2}.
\]
固定 $\Delta_1$, 取 $\Delta>\max\left\{\Delta_1,\ln\frac{2|\mathrm{e}^{\Delta_1}f(\Delta_1)|}{\varepsilon}\right\}$, 当 $x>\Delta$ 时, 有
\[
\begin{aligned}
|f(x)-0|&=\left|\frac{\mathrm{e}^{x}f(x)}{\mathrm{e}^{x}}\right|\\
&=\left|\frac{\mathrm{e}^{x}f(x)-\mathrm{e}^{\Delta_1}f(\Delta_1)}{\mathrm{e}^{x}-\mathrm{e}^{\Delta_1}}\cdot\frac{\mathrm{e}^{x}-\mathrm{e}^{\Delta_1}}{\mathrm{e}^{x}}+\frac{\mathrm{e}^{\Delta_1}f(\Delta_1)}{\mathrm{e}^{x}}\right|
\end{aligned}
\]
\[
\begin{aligned}
&\xlongequal[\exists\xi\in(\Delta_1,x)]{\text{Cauchy 中值定理}}\left|\frac{\mathrm{e}^{\xi}[f(\xi)+f'(\xi)]}{\mathrm{e}^{\xi}}(1-\mathrm{e}^{\Delta_1-x})+\frac{\mathrm{e}^{\Delta_1}f(\Delta_1)}{\mathrm{e}^{x}}\right|\\
&\leqslant|f(\xi)+f'(\xi)|+\frac{|\mathrm{e}^{\Delta_1}f(\Delta_1)|}{\mathrm{e}^{x}}<\frac{\varepsilon}{2}+\frac{|\mathrm{e}^{\Delta_1}f(\Delta_1)|}{\mathrm{e}^{\Delta}}\\
&<\frac{\varepsilon}{2}+\frac{\varepsilon}{2}=\varepsilon,
\end{aligned}
\]
所以,
\[
\lim_{x\to+\infty}f(x)=0.
\]
\end{proof}

\begin{exblock}{练习题 3.4}
  \item 求下列不定型的极限:
  \begin{enumerate}[label=(\arabic*)]
    \item $\lim\limits_{x\to0}\frac{\mathrm{e}^{x}-1}{\sin x}$;
    \item $\lim\limits_{x\to\frac{\pi}{6}}\frac{1-2\sin x}{\cos 3x}$;
    \item $\lim\limits_{x\to0}\frac{\ln(1+x)-x}{\cos x-1}$;
    \item $\lim\limits_{x\to0}\frac{\tan x-x}{x-\sin x}$;
    \item $\lim\limits_{x\to\frac{\pi}{2}}\frac{\tan x-6}{\sec x+5}$;
    \item $\lim\limits_{x\to0}\left(\frac{1}{x}-\frac{1}{\mathrm{e}^{x}-1}\right)$;
    \item $\lim\limits_{x\to0}(\tan x)^{\sin x}$;
    \item $\lim\limits_{x\to1}x^{\frac{1}{1-x}}$;
    \item $\lim\limits_{x\to0}(1+x^{2})^{\frac{1}{x}}$;
    \item $\lim\limits_{x\to0^{+}}\sin x\ln x$;
    \item $\lim\limits_{x\to0}\left(\frac{1}{x^{2}}-\frac{1}{\sin^{2}x}\right)$;
    \item $\lim\limits_{x\to0}\left(\frac{\tan x}{x}\right)^{\frac{1}{x^{2}}}$;
    \item $\lim\limits_{x\to0}\frac{\mathrm{e}^{ax}-\mathrm{e}^{bx}}{\sin ax-\sin bx}$;
    \item $\lim\limits_{x\to0}\frac{x\cot x-1}{x^{2}}$;
    \item $\lim\limits_{x\to0}\frac{x(\mathrm{e}^{x}+1)-2(\mathrm{e}^{x}-1)}{x^{3}}$;
    \item $\lim\limits_{x\to+\infty}\frac{\mathrm{e}^{x}+\mathrm{e}^{-x}}{x^{3}}$;
    \item $\lim\limits_{x\to0^{+}}x^{x}$;
    \item $\lim\limits_{x\to1^{-}}\ln x\ln(1-x)$;
    \item $\lim\limits_{x\to+\infty}\left(\cos\frac{a}{x}\right)^{x}$;
    \item $\lim\limits_{x\to+\infty}\left(\frac{2}{\pi}\arctan x\right)^{x}$;
    \item $\lim\limits_{x\to\infty}x\left(\left(1+\frac{1}{x}\right)^{x}-\mathrm{e}\right)$;
    \item $\lim\limits_{x\to\infty}\left(\frac{2}{\pi}\arccos x\right)^{\frac{1}{x}}$;
    \item $\lim\limits_{x\to\infty}(\pi-2\arctan x)\ln x$;
    \item $\lim\limits_{x\to+\infty}\left(\frac{\pi}{2}-\arctan x\right)^{\frac{1}{\ln x}}$;
    \item $\lim\limits_{x\to0}\left(\frac{\ln(1+x)^{1+x}}{x^{2}}-\frac{1}{x}\right)$.
  \end{enumerate}

  \item 设 $f$ 在 $0$ 的某开邻域内 2 阶连续可导且 $f(0)=0$, 定义
  \[
  g(x)=\begin{cases}\frac{f(x)}{x}, & x\neq0,\\ f'(0), & x=0.\end{cases}
\]
  证明: $g$ 在这开邻域内是连续可导的.

  \item (1) 设函数 $f$ 在 $x_0$ 处具有连续的 2 阶导数. 两次应用 L'Hospital 法则证明:
  \[
  \lim_{h\to0}\frac{f(x_0+h)+f(x_0-h)-2f(x_0)}{h^{2}}=f''(x_0);
\]
  (2) 当函数 $f$ 在点 $x_0$ 处具有 2 阶导数时, 上述等式是否仍成立.

  \item 设 $f(0)=0$, $f'$ 在原点的某开邻域内连续, 且 $f'(0)\neq0$. 证明: $\lim\limits_{x\to0^{+}}x^{f(x)}=1$.

  \item 证明: $f(x)=x^{3}\mathrm{e}^{-x^{2}}$ 为有界函数.
\end{exblock}

\begin{exblock}{思考题 3.4}
  \item 设 $x_1=\sin x_0>0$, $x_{n+1}=\sin x_n$, $x\in\mathbb{N}$. 证明:
  \[
  \lim_{n\to+\infty}\sqrt{\frac{n}{3}}x_n=1.
\]

  \item 应用 L'Hospital 法则或数学归纳法证明:
  \[
  \lim_{x\to0}\frac{1-\cos a_1x\cdots\cos a_nx}{x^{2}}=\frac{1}{2}\sum_{k=1}^{n}a_k^{2}.
\]
\end{exblock}
